\section{Related work and open questions}\label{sec:related}
\subsection{Related work}
\emph{Arithmetic matrices.} The fixed-index part of the connection in
Section~\ref{sec:finite} follows a pattern used by Hilberdink~\cite{Hil17}
for multiplicative Toeplitz matrices $(f(i/j)\one_{\{j\mid i\}})$: after
normalization, Gram matrices of the finite matrices converge to a positive
Hilbert--Schmidt operator, so each fixed singular value has a limit
described by an eigenvalue of that operator. Hilberdink and
Pushnitski~\cite{HP21,HP24} then determined eigenvalue asymptotics for
infinite arithmetic matrices built from least common multiples, such as
$(n^\sigma m^\sigma/[n,m]^\tau)$, obtaining power laws such as
$\lambda_n\sim\kappa n^{-\rho}$. Their proofs use a factorization of these
matrices as infinite tensor products over primes, which expresses the
eigenvalues as products of local eigenvalues, followed by a Tauberian
theorem of Wiener--Ikehara type. They also show that the constant $\kappa$
need not be the one suggested by the diagonal entries; see~(1.5) of
\cite{HP21} and the discussion following it. The kernel $K$ is of a
different kind: its entries depend
on the ancestor relation in the tree, not on a multiplicative function of
the indices, and we have not found a tensor-product factorization of it.
Its eigenvalue law also carries a logarithmic factor rather than being a
pure power. In contrast with the example of~\cite{HP21}, for $K$ the
leading constant is the one given by the count of diagonal masses.

\emph{Operators on trees.} The synthesis map $V$ of
\eqref{eq:gram-identification} is a discrete Hardy-type operator: for
finitely supported $x$, $(Vx)(\omega)$ is the sum of $x_d$ over the
vertices $d$ on the path determined by $\omega$. Evans, Harris and
Lang~\cite{EHL} proved, for Hardy operators on metric trees with finite
vertex degrees and suitable weights $u,v$, that the approximation numbers
satisfy $n\,a_n\to\alpha_p\int|uv|$, extending earlier Hilbert-space
results cited there. When this limit is finite,
$a_n=O(1/n)$. For $V$, Theorem~\ref{thm:main} and
\eqref{eq:gram-identification} give
$a_n(V)=(\lambda_n(K)/c_0)^{1/2}\sim(n\log n)^{-1/2}$, so $V$ is not covered
by that theorem with a finite integral. The tree of
Section~\ref{sec:kernel} has vertices of infinite degree, but this is a
feature of the presentation rather than of the operator. Revealing the
primes one at a time gives a binary filtration of $\Omega$ whose cylinders
$B$ are determined by the selected primes up to some prime $p$. Each $A_d$
is the cylinder of this kind with $p=P^+(d)$ in which the last revealed
prime is selected, and
\[
 T=\sum_B c_B\,\nu(B)^{-1}\rankone{\one_B},
 \qquad
 c_B=\begin{cases}\nu(B),&B=\Omega\text{ or the last revealed prime of }B
 \text{ is selected},\\0,&\text{otherwise}.\end{cases}
\]
Results on weighted dyadic or martingale Carleson embeddings, and on
singular values of weighted operators on trees of bounded degree, are
therefore relevant comparisons. We have not determined whether any of them
yields the exact constant in Theorem~\ref{thm:main}, only two-sided
bounds, or neither.

\emph{Ultrametric spaces.} For almost every pair $(\omega,\omega')$, the
kernel of $T$ is one plus the length of the longest common initial segment
of the increasing sequences of primes selected by $\omega$ and $\omega'$,
equivalently, the number of vertices whose cylinders contain both points.
This is an ultrametric quantity. Bendikov, Grigor'yan, Pittet and
Woess~\cite{BGPW} studied isotropic Laplacians on ultrametric spaces with
compact balls and a reference measure, which we again call $\nu$; their
eigenfunctions are the functions $\one_C/\nu(C)-\one_B/\nu(B)$, for a ball
$C$ with parent ball $B$~\cite[Theorem~3.8]{BGPW}. The operator $T$ is not diagonal in the
analogous basis. For $h=3\one_{A_2}-1$, which is this function for the
cylinder $A_2$ and its parent $\Omega$, one has $\inner{\one}{h}=0$ and
$(Th)(\omega)=2\sum\nu_d$ for $\omega\in A_2$, the sum taken over the
vertices $d\ne1$ on the path of $\omega$. This is not constant on $A_2$:
on $A_6$ it is at least $2\nu_2+2\nu_6=2/3+1/6$, while on $A_{10}$ the
vertices after $2$ have total mass at most $\nu_{10}/(1-1/7)=7/144<1/12$,
because every later child ratio is at most $1/7$ by \eqref{eq:arith-ratios};
so $Th<2/3+1/6$ there. Hierarchical Schr\"odinger operators, which add a potential
to a hierarchical Laplacian, are studied by Bendikov, Grigor'yan and
Molchanov~\cite{BGM}, including questions about discrete spectrum. Their
stated setting is homogeneous, and we have not carried out a comparison
with the inhomogeneous operator $T$.

\emph{Self-similar spectral problems.} The recursion behind
Proposition~\ref{prop:sandwich}, in which a subtree operator is a rank-one
operator plus rescaled copies of smaller subtree operators, resembles the
recursive counting identities used for spectral asymptotics on
self-similar sets, where renewal theory supplies the asymptotics; see
Kigami and Lapidus~\cite{KL}. Appendix~\ref{sec:stationary} uses the key
renewal theorem in exactly this way for stationary trees. For the
arithmetic tree the subtrees are not rescaled copies of a single operator,
and the asymptotics come instead from comparison with vertex counts.

\emph{The arithmetic count.} Theorem~\ref{thm:arithmetic} is an
application of classical estimates for smooth numbers~\cite{HT}, of the
same kind as those used for sums involving the largest prime factor; for
instance, Erd\H{o}s, Ivi\'c and Pomerance~\cite{EIP} evaluate
$\sum_{n\le x}1/P^+(n)$ through $\Psi(x,y)$ and $\rho$. We do not regard
the count as new in method.

We have not found Theorem~\ref{thm:main} in the sources discussed above.
The comparisons described there as not carried out could change that
assessment.

\subsection{Open questions}
\begin{question}\label{q:drop}
Under \eqref{eq:tree-assumptions}, suppose $D$ is regularly varying at
zero with index $-1$. Must $\ncount(T;t)\sim D(t)$ hold without either
vanishing child ratios or the profile hypotheses of
Theorem~\ref{thm:profile}?
\end{question}

The exact decomposition \eqref{eq:exact-shifts} reduces this question to
the locations of unit-integral, zero--one profiles at varying mass scales.
Their common integral does not, by itself, justify exchanging a limit
and the vertex sum. Proposition~\ref{prop:integrated} gives an integrated
comparison without profile hypotheses, but does not establish a pointwise
counting asymptotic. Proposition~\ref{prop:kary} is not a counterexample to
Question~\ref{q:drop}, because it lacks the stipulated regular variation.

Trees whose splitting rule depends on the position of a vertex's mass on
a logarithmic scale are natural candidates for a negative answer. For such
trees we have proved neither that the mass count is regularly varying nor
that $\ncount(T;t)/D(t)$ fails to converge, and we make no claim about
them here.

\begin{question}
Can the proof of Theorem~\ref{thm:arithmetic} and the sandwich in
Proposition~\ref{prop:sandwich} be made quantitative enough to give a
useful remainder in \eqref{eq:maincount}?
\end{question}
The proof as written takes $X\to\infty$ for fixed $a,b$, then sends
$a\downarrow0$ and $b\uparrow1$. The compact-$u$ input and this order
of limits yield a leading equivalent only. In particular, the present
argument does not identify a second term or a rate.

\begin{question}
Can one bound $\norm{C_n^*C_n/n-K}$ explicitly in a range that makes
Corollary~\ref{cor:growing} effective for useful $r_n$?
\end{question}
Hilbert--Schmidt convergence suffices for fixed indices, but supplies no
explicit scale in the proof given here. A quantitative treatment must
control both the rough-number densities and the matrix tail, then compare
that error with $1/(r_n\log r_n)$.

\begin{remark}[Arithmetic reach]
The main result concerns a positive kernel whose signs can be removed by
a diagonal unitary. Its proof uses classical smooth-number estimates and
does not produce new cancellation in M\"obius sums. An application to a
zero-free region would require an additional theorem that controls the
relevant signed arithmetic quantity; the asymptotic
$\lambda_r(K)\sim c_0/(r\log r)$ is not such a theorem by itself.
\end{remark}

\begin{remark}[Attribution]
The compactness argument, interlacing, the rank-one trace identity, the
renewal theorem, and the smooth-number inputs belong to established
operator theory, probability, and analytic number theory. We make no
priority claim for the abstract tree results in Section~\ref{sec:tree}
or in the appendices.
\end{remark}

\subsection*{Role of AI}
AI assisted the research development, literature search, proof exposition,
and mathematical checking under the author's direction. Process-separated
AI reviews are not independent expert peer review, and no independent human
expert has examined the proofs. The release repository records the
inspected sources, the proof reviews, and their dispositions.
